% % % % % % % % % % % % % % % % % % % % % % % % % % % % % % % % % % % % % % % % % % % % % % % % % % % % % % % % % % % % % % % %
% SAT_STYLE_MATH_MCQS.tex
% 1_07.tex
% 
% encoding: UTF-8 
% EOL: LF
%
% format: LaTeX
% indent: spaces (2)
% width: 127
% % % % % % % % % % % % % % % % % % % % % % % % % % % % % % % % % % % % % % % % % % % % % % % % % % % % % % % % % % % % % % % %
Ici, $111\dots 11$ désigne l'entier écrit (seulement) avec $N$ chiffres $1$. La somme $1 + 11 + 111 + \cdots + 111\dots 11$ %
vaut:
%
\begin{enumerate}
  \item $\frac{1}{9}\!\left(\frac{10^{N+1}-1}{9} - (N+1)\right)$
  \item $\frac{1}{9}\!\left(\frac{10^{N+2}-1}{9} - (N+1)\right)$
  \item $\frac{1}{9}\!\left(\frac{10^{N}-1}{9} - (N+1)\right)$
  \item $\frac{1}{9}\!\left(\frac{10^{N}-1}{9} - N\right)$
  \item $\frac{10^{N+1}-(N+1)}{81}$
\end{enumerate}
%
\begin{proof}
Pour établir la somme $S(N) = 1 + 11 + 111 + \cdots + 111\dots 11$\;\,($\naturals{N}$), notons que  
%
\begin{align}
  9 \cdot S(N) & = 9 + 99 + \cdots + \overbrace{9 \dots 9}^{N \text{ `$9$'}} \\
  & = (10-1) + (10^2 - 1) + \cdots + (10^N - 1) \\
  & = (10 + 10^2 + \cdots + 10^N) - N \\
  & = (1 + 10 + 10^2 + \cdots + 10^N) - (N + 1)
\end{align}
%
L'application à $q=10$ de la formule 
%
\begin{equation}
  q^{N+1} - 1  = (q-1)(1 + \cdots + q^N)
\end{equation} 
%
donne 
% 
\begin{equation}
  9 \cdot S(N) = \frac{10^{N+1} - 1}{9} - (N+1).
\end{equation}
%
Réponse \textbf{A}.
\end{proof}